\section{Motivation}\label{sec:motivation}

Consider evaluating the diversity of outputs from a language model policy $\pi$.
This is useful for comparing any generation process: RLHF-trained or finetuned models, different prompting strategies, decoding parameters, representation-engineering interventions, or even human writers.
Existing approaches typically measure diversity via embedding-based metrics \citep{du2019boosting} or surface-level $n$-gram statistics.
A complementary line of work applies information theory directly \citep{li2016diversity, zhang2018aim}, but uses pairwise mutual information between input and response as a training or decoding objective rather than as an evaluation-time diagnostic (contrasted in Section~\ref{sec:pairwise-mi}).
These existing approaches have known limitations: embedding distances may not align with semantically meaningful differences, and $n$-gram metrics conflate lexical variation with genuine content diversity.

We propose an alternative grounded in information theory: use a trusted base model $\theta$ with per-token logprob access to assess how much information responses from $\pi$ share with one another.
The core intuition is that if $\pi$'s responses are diverse, then seeing one response should not help $\theta$ predict another.
If $\pi$'s responses are repetitive or constrained to a narrow set of modes, then conditioning on one response should significantly reduce $\theta$'s surprise at subsequent responses.

\paragraph{Why a separate base model?} We deliberately avoid using $\pi$ to judge its own outputs.
A post-trained model may assign high probability to all its (narrowed) outputs, masking the loss of diversity.
Using an independent $\theta$ provides an external reference frame.

\paragraph{Important caveat.} The entire framework rests on $\theta$'s ability to perform in-context learning: when conditioned on previous responses, $\theta$ must be able to recognize patterns and reduce its surprise at subsequent similar responses.
All quantities derived below (the progressive conditional surprise curve, the asymptotic floor, and the coherence term) measure properties of $\pi$'s outputs \emph{as perceived by $\theta$}.
If $\pi$'s outputs differ only in ways that $\theta$ cannot distinguish from context, the metric underestimates diversity.
Conversely, if $\theta$ hallucinates structure that is not present, it overestimates redundancy.
The metric therefore measures diversity \emph{relative to the base model's perceptual capabilities}, which is a weaker claim than measuring the true diversity of $\pi$.
As base models continue to improve at in-context learning, this gap narrows.

\paragraph{Contributions.}
\begin{itemize}[leftmargin=*, itemsep=2pt]
  \item We define a \emph{progressive conditional surprise curve} $a_k$ for language-model outputs, computed by a single forward pass of a base model $\theta$ over the concatenated responses, whose shape characterizes the learnable structure among outputs (Section~\ref{sec:progressive-conditioning}).
  \item We introduce a separate \emph{coherence} term $C$ measuring per-byte plausibility of outputs under $\theta$, and combine it with the asymptotic curve value $a_\infty$ to yield a single plausibility-weighted diversity score $\mathbf{D_{Ca_n} = C \times a_n}$ (Sections~\ref{sec:coherence} and \ref{sec:scalar}).
  \item We validate the metric on synthetic mode scenarios, on Tevet and Berant's diversity-eval benchmark (where $C \times a_n$ is competitive with embedding-based metrics), on a cross-model scaling study, and on the OLMo-2-7B post-training pipeline (Section~\ref{sec:experiments}).
\end{itemize}

\paragraph{Evidence preview.}
On Tevet and Berant's binary model-comparison test, $C \times a_n$ reaches ROC AUC $= \tevetMcDivPromptGenAUC{}$ on McDiv prompt\_gen, competitive with sentence-embedding metrics (Section~\ref{sec:tevet}).
In synthetic mode-count experiments, $C \times a_n$ (aggregated over \modeCountNDraws{} draws per mode count) increases monotonically with mode count $m$ for both GPT-2 (Pearson $r = \modeCountCxAnPearsonGPT{}$) and Qwen2.5-3B ($r = \modeCountCxAnPearsonQwen{}$) (Section~\ref{sec:mode-count-scaling}).
A cross-model scaling study across \scalingNumLlamaModels{} Llama sizes and an OLMo-2-7B post-training pipeline (base~$\to$~SFT~$\to$~DPO~$\to$~RLVR) provide further corroboration.

\paragraph{Who should care.}
The metric applies to anyone comparing output diversity across generation processes (post-trained models, prompting or decoding strategies, representation-engineering interventions, or human writers) given a base model $\theta$ with per-token logprob access.
It requires only a single forward pass of $\theta$ over the responses; no trained embedding model, reference corpus, or human labels are needed.
Practitioners diagnosing post-training diversity collapse (e.g., in RLHF or DPO) can use $C \times a_n$ alongside existing embedding- or $n$-gram-based diagnostics.
Researchers evaluating model-steering interventions (activation steering, representation engineering, SAE-based feature steering) can similarly quantify how a steering strength trades off against output diversity.
Because the metric measures diversity relative to $\theta$'s in-context learning capability, it also strengthens automatically as base models improve.

Figure~\ref{fig:pipeline} sketches the full pipeline.
The remainder of the paper formalizes each step and validates the resulting metric.

\begin{figure}[tbp]
\centering
\begin{tikzpicture}[
    font=\footnotesize,
    box/.style={draw, rounded corners=2pt, align=center, inner sep=3pt},
    stage/.style={box, fill=gray!8},
    cond/.style={box, fill=orange!10},
    uncond/.style={box, fill=blue!8},
    final/.style={box, fill=green!10, very thick},
    arrow/.style={-{Latex[length=1.8mm]}, thick},
    node distance=2mm and 4mm,
]

% --- INPUT + FORMAT (combined, top) ---
\node[stage, text width=0.92\textwidth] (input) {
    \textbf{Input:} prompt $p$, responses $r_1,\ldots,r_n$ from policy $\pi$.\\[2pt]
    \textbf{Format:} \texttt{p\textbackslash n\textbackslash nResponse A: $r_{\sigma(1)}$\textbackslash n\textbackslash nResponse B: $r_{\sigma(2)}$\,\ldots}\\[2pt]
    \textbf{Tokenize:} $[\texttt{Resp}][\texttt{onse}][\texttt{ A}][\texttt{:}][\texttt{ Rain}][\texttt{ falls}]\cdots$
};

% --- CONDITIONAL TRACK (LEFT) ---
\node[cond, below left=5mm and 1mm of input.south, text width=0.45\textwidth, anchor=north east] (cond) {
    \textbf{Conditional track} \;{\scriptsize ($\times\,n_\sigma$ permutations)}\\[2pt]
    Feed concatenated context to base model $\theta$ in one forward pass.\\
    For each response slot $k$:
    \[ a_k^\sigma \;=\; -\!\!\sum_{t \in r_{\sigma(k)}}\!\!\log_2 \theta(t \mid \mathrm{prev}) \]
    Divide by $\|r_{\sigma(k)}\|$ \emph{(per-byte)}.\\
    Average over $\sigma$ to get $\bar{a}_k$ curve.\\[2pt]
    Last point: $a_n \approx a_\infty$
};

% --- UNCONDITIONAL TRACK (RIGHT) ---
\node[uncond, below right=5mm and 1mm of input.south, text width=0.45\textwidth, anchor=north west] (unc) {
    \textbf{Unconditional track} \;{\scriptsize ($n$ separate passes)}\\[2pt]
    Feed each $r_i$ alone (with prompt $p$) to $\theta$.
    \[ h(r_i\mid p) \;=\; -\tfrac{1}{\|r_i\|}\log_2 \theta(r_i \mid p) \]
    \emph{(per-byte)} Take the geometric mean:
    \[ C \;=\; 2^{-\,\frac{1}{n}\sum_i h(r_i\mid p)} \]
};

% Arrows from input down to each track
\draw[arrow] (input.south) -- ++(0,-1mm) -| (cond.north);
\draw[arrow] (input.south) -- ++(0,-1mm) -| (unc.north);

% --- FINAL FORMULA ---
\node[final, below=4mm of input, text width=0.62\textwidth, yshift=-49mm] (out) {
    \textbf{Diversity score:}\quad $D_{Ca_n} \;=\; C \times a_n$\quad{\scriptsize(bits/byte)}
};

% Arrows from each track straight down into the final box
\draw[arrow] (cond.south) -- (cond.south |- out.north);
\draw[arrow] (unc.south)  -- (unc.south |- out.north);

\end{tikzpicture}
\caption{The diversity-metric pipeline.
Given prompt $p$ and $n$ responses from policy $\pi$, we format them with response labels (``Response A:'', ``Response B:'', \ldots) and tokenize.
The \textcolor{orange!70!black}{\textbf{conditional track}} (left) feeds the concatenated context to the base model $\theta$ in a single forward pass per permutation $\sigma$, extracts per-response total surprise, divides by each response's UTF-8 byte count, and averages over permutations to get $\bar a_k$.
Its last point $a_n$ estimates the asymptotic floor $a_\infty$.
The \textcolor{blue!70!black}{\textbf{unconditional track}} (right) scores each $r_i$ independently against the prompt alone, again per-byte; the geometric mean of these surprises defines coherence $C = 1/\mathrm{PPL}_\theta(\pi, p)$, the reciprocal of the geometric-mean per-byte perplexity.
The product $D_{Ca_n} = C \times a_n$ measures plausibility-weighted residual diversity (bits/byte).
Formal definitions: Sections~\ref{sec:setup-notation}, \ref{sec:progressive-conditioning}, \ref{sec:coherence}, \ref{sec:scalar}.}
\label{fig:pipeline}
\end{figure}

